\section{Additional conclusions recorded in the introduction}

Besides the main theorems above, the introduction records several qualitative conclusions about strong solutions.

\begin{enumerate}[leftmargin=2em]
    \item Strong solutions satisfy a comparison principle.
    \item Strong solutions preserve the zero set in time.
    \item Distinct intervals of positivity do not interact with one another.
    \item The analysis therefore reduces to positive strong solutions on the three model intervals
    \[
        \R, \qquad \R_+, \qquad (0,L).
    \]
\end{enumerate}

The introduction also emphasizes two kinds of pathology that must be ruled out before time $\kappa^{-1}$:
\begin{enumerate}[leftmargin=2em]
    \item blow-up at infinity;
    \item blow-up at zero.
\end{enumerate}

Finally, the paper relates this PDE to a forward-backward stochastic differential equation. In that interpretation, the solution $u$ is a decoupling function, and the root-Lipschitz control of $\sqrt{u}$ is one of the key quantities required in the companion stochastic analysis.
